% !TeX root = ../main.tex
% ---------------------------------------------------------------------------
% Chapter imported from Weekly Study Notes/ Chapter~\ref{ch:classification}.tex
% Initial enrichment pass: copied directly from the weekly study-note body.
% ---------------------------------------------------------------------------
% ======================================================================
\chapter{Rock Mass Classification: RMR, Q-System and GSI}
\label{ch:classification}
% ======================================================================

\section{Purpose and workflow of classification}
\label{sec:purpose-and-workflow-of-classification}

\subsection{Rock Mass Classification -- Overview}
% ======================================================================

Rock mass classification systems were developed to:
\begin{itemize}
 \item identify the most significant parameters influencing rock mass behaviour;
 \item divide rock masses into groups of similar engineering behaviour;
 \item provide quantitative data and guidelines for preliminary design;
 \item support communication between engineers and geologists; and
 \item provide input to empirical support systems and rock mass failure criteria.
\end{itemize}

Key contributors and systems include:
\begin{itemize}
 \item Terzaghi (1946): Rock Load Factor;
 \item Bieniawski (1973, 1989): Rock Mass Rating (RMR);
 \item Barton et al.\ (1974): Q-system;
 \item Hoek et al.\ (1994--1995): Geological Strength Index (GSI);
 \item Molinda and Mark (1994): Coal Mine Roof Rating (CMRR).
\end{itemize}

\begin{keybox}[Classification is not calculation only]
Rock mass classification requires both numerical calculation and geological judgement.
Engineers must justify the selected ratings, not only report final RMR, Q, or GSI values.
\end{keybox}

% ======================================================================
\subsection{Classification Workflow}
% ======================================================================

The recommended workflow Chapter~\ref{ch:classification} is shown in Figure~\ref{fig:classification_workflow}.

\begin{figure}[H]
\centering
\resizebox{\textwidth}{!}{%
\begin{tikzpicture}[node distance=0.4cm, font=\small]

\node[smallworkflowbox=orange] (data)
{\textbf{1. Collect data}\\
\footnotesize UCS, RQD, joints, water, stress};

\node[smallworkflowbox=green, right=of data] (rmr)
{\textbf{2. Calculate RMR}\\
\footnotesize \(R_1+R_2+R_3+R_4+R_5\)};

\node[smallworkflowbox=blue, right=of rmr] (q)
{\textbf{3. Calculate Q}\\
\footnotesize \(RQD/J_n \times J_r/J_a \times J_w/SRF\)};

\node[smallworkflowbox=purple, right=of q] (gsi)
{\textbf{4. Estimate GSI}\\
\footnotesize chart + correlations};

\node[smallworkflowbox=teal, right=of gsi] (check)
{\textbf{5. Compare}\\
\footnotesize justify final design value};

\draw[workflowarrow] (data) -- (rmr);
\draw[workflowarrow] (rmr) -- (q);
\draw[workflowarrow] (q) -- (gsi);
\draw[workflowarrow] (gsi) -- (check);

\draw[feedbackarrow] (check.south) -- ++(0,-0.6) -| node[pos=0.25,below,font=\footnotesize]
{revise if classifications are inconsistent} (data.south);

\end{tikzpicture}
}
\caption{Recommended rock mass classification workflow.}
\label{fig:classification_workflow}
\end{figure}

\begin{warnbox}[Do not force all systems to agree]
RMR, Q, and GSI measure related but not identical aspects of rock mass quality.
Differences between them are not necessarily mistakes. A good report explains why
the values differ and selects a defensible design value.
\end{warnbox}

% ======================================================================
\section{Rock Mass Rating (RMR) System}
% ======================================================================

The RMR system of Bieniawski (1989) is based on five parameters:
\[
\text{RMR} = R_1 + R_2 + R_3 + R_4 + R_5
\]
where:
\begin{itemize}
 \item \(R_1\) = intact rock strength rating;
 \item \(R_2\) = RQD rating;
 \item \(R_3\) = discontinuity spacing rating;
 \item \(R_4\) = discontinuity condition rating;
 \item \(R_5\) = groundwater rating.
\end{itemize}

\begin{notebox}[Basic RMR vs adjusted RMR]
In Chapter~\ref{ch:classification}, calculate the \textbf{basic RMR} from the five parameters above.
Orientation adjustment is used later for design applications such as tunnels and foundations. For GSI conversion, use \(RMR_{\text{basic}}\), not an orientation-adjusted RMR.
\end{notebox}

% ----------------------------------------------------------------------
\paragraph{R1: Strength of Intact Rock Material}

This rating is based on uniaxial compressive strength or point load index.

\begin{table}[H]
\centering
\small
\begin{tabularx}{0.85\textwidth}{l l c}
\toprule
\textbf{Point Load \(I_{s(50)}\) (MPa)} & \textbf{UCS (MPa)} & \textbf{Rating} \\
\midrule
\(>10\) & \(>250\) & 15 \\
4--10 & 100--250 & 12 \\
2--4 & 50--100 & 7 \\
1--2 & 25--50 & 4 \\
-- & 5--25 & 2 \\
-- & 1--5 & 1 \\
-- & \(<1\) & 0 \\
\bottomrule
\end{tabularx}
\caption{RMR rating for intact rock strength.}
\label{tab:rmr_strength}
\end{table}

% ----------------------------------------------------------------------
\paragraph{R2: Rock Quality Designation (RQD)}

\[
\text{RQD} =
\frac{
\sum \text{length of sound core pieces} > 100~\text{mm}
}
{\text{total core run length}}
\times 100\%
\]

\begin{table}[H]
\centering
\small
\begin{tabular}{c c l}
\toprule
\textbf{RQD (\%)} & \textbf{Rating} & \textbf{Description} \\
\midrule
90--100 & 20 & Excellent \\
75--90 & 17 & Good \\
50--75 & 13 & Fair \\
25--50 & 8 & Poor \\
\(<25\) & 3 & Very poor \\
\bottomrule
\end{tabular}
\caption{RMR rating for RQD.}
\label{tab:rmr_rqd}
\end{table}

\begin{warnbox}[RQD is not strength]
RQD is a geometric index. It describes the percentage of core pieces longer than
100 mm. It does not directly measure joint roughness, joint infilling, shear strength,
or deformability.
\end{warnbox}

% ----------------------------------------------------------------------
\paragraph{R3: Spacing of Discontinuities}

\begin{table}[H]
\centering
\small
\begin{tabular}{c c}
\toprule
\textbf{Joint spacing (m)} & \textbf{Rating} \\
\midrule
\(>2\) & 20 \\
0.6--2 & 15 \\
0.2--0.6 & 10 \\
0.06--0.2 & 8 \\
\(<0.06\) & 5 \\
\bottomrule
\end{tabular}
\caption{RMR rating for discontinuity spacing.}
\label{tab:rmr_spacing}
\end{table}

If joint spacing is not directly measured, it may be estimated from RQD using the
Priest--Hudson relationship:
\[
RQD = 110.4 - 3.68\lambda
\qquad
(6 < \lambda < 16)
\]
where \(\lambda\) is the joint frequency in joints per metre.

The average joint spacing is then:
\[
s = \frac{1}{\lambda}
\]

\begin{calcbox}[Example calculation]
If \(RQD = 68\%\), then:
\[
\lambda = \frac{110.4 - 68}{3.68} = 11.52~\text{joints/m}
\]
\[
s = \frac{1}{11.52}=0.087~\text{m}
\]
This spacing would fall in the 0.06--0.2 m category.
\end{calcbox}

\begin{warnbox}[Priest--Hudson validity range]
The linear approximation \(RQD = 110.4 - 3.68\lambda\) is only valid for
\(6 < \lambda < 16\), approximately corresponding to RQD values between 51\% and 88\%.
For very poor rock, use this only as a rough guide and explain the uncertainty.
\end{warnbox}

% ----------------------------------------------------------------------
\paragraph{R4: Condition of Discontinuities}

\begin{table}[H]
\centering
\small
\begin{tabularx}{\textwidth}{Y c}
\toprule
\textbf{Condition description} & \textbf{Rating} \\
\midrule
Very rough surfaces, not continuous, no separation, unweathered wall rock & 30 \\
Slightly rough surfaces, separation \(<1\) mm, slightly weathered walls & 25 \\
Slightly rough surfaces, separation \(<1\) mm, highly weathered walls & 20 \\
Slickensided surfaces OR gouge \(<5\) mm thick OR separation 1--5 mm, continuous & 10 \\
Soft gouge \(>5\) mm thick OR separation \(>5\) mm, continuous & 0 \\
\bottomrule
\end{tabularx}
\caption{RMR rating for discontinuity condition.}
\label{tab:rmr_condition}
\end{table}

\begin{notebox}[Use the most representative condition]
If the joint description contains mixed features, choose the rating that best represents
the likely engineering behaviour and justify the decision. Do not simply choose the most
favourable description.
\end{notebox}

% ----------------------------------------------------------------------
\paragraph{R5: Groundwater Conditions}

\begin{table}[H]
\centering
\small
\begin{tabularx}{\textwidth}{l l l c}
\toprule
\textbf{Inflow per 10 m tunnel (L/min)} & \textbf{\(u/\sigma_1\)} & \textbf{General condition} & \textbf{Rating} \\
\midrule
None & 0 & Completely dry & 15 \\
\(<10\) & 0--0.1 & Damp & 10 \\
10--25 & 0.1--0.2 & Wet & 7 \\
25--125 & 0.2--0.5 & Dripping & 4 \\
\(>125\) & \(>0.5\) & Flowing & 0 \\
\bottomrule
\end{tabularx}
\caption{RMR rating for groundwater.}
\label{tab:rmr_groundwater}
\end{table}

% ----------------------------------------------------------------------
\paragraph{Rock Mass Classes Based on RMR}

\begin{table}[H]
\centering
\small
\begin{tabular}{c c l c c}
\toprule
\textbf{RMR} & \textbf{Class} & \textbf{Description} & \textbf{Cohesion (kPa)} & \textbf{Friction angle} \\
\midrule
81--100 & I & Excellent & \(>400\) & \(>45^\circ\) \\
61--80 & II & Good & 300--400 & 35--45\(^{\circ}\) \\
41--60 & III & Fair & 200--300 & 25--35\(^{\circ}\) \\
21--40 & IV & Poor & 100--200 & 15--25\(^{\circ}\) \\
\(<20\) & V & Very poor & \(<100\) & \(<15^\circ\) \\
\bottomrule
\end{tabular}
\caption{Rock mass classes based on RMR.}
\label{tab:rmr_classes}
\end{table}

% ======================================================================
\section{NGI Q-System}
% ======================================================================

The Q-system of Barton et al.\ (1974) is defined as:
\[
Q =
\frac{RQD}{J_n}
\times
\frac{J_r}{J_a}
\times
\frac{J_w}{SRF}
\]

The three terms have physical meaning:
\[
\frac{RQD}{J_n} = \text{block size}
\]
\[
\frac{J_r}{J_a} = \text{inter-block shear strength}
\]
\[
\frac{J_w}{SRF} = \text{active stress and water effect}
\]

\begin{notebox}[RQD is entered as a percentage]
Use \(RQD=70\), not \(0.70\). A common error is to enter the decimal form, which
makes the Q-value 100 times too small.
\end{notebox}

% ----------------------------------------------------------------------
\paragraph{Joint Set Number \(J_n\)}

\begin{table}[H]
\centering
\small
\begin{tabularx}{0.85\textwidth}{Y c}
\toprule
\textbf{Joint sets} & \textbf{\(J_n\)} \\
\midrule
Massive, no or few joints & 0.5--1 \\
One joint set & 2 \\
One joint set plus random joints & 3 \\
Two joint sets & 4 \\
Two joint sets plus random joints & 6 \\
Three joint sets & 9 \\
Three joint sets plus random joints & 12 \\
Four or more joint sets, heavily jointed & 15 \\
Crushed rock, earthlike & 20 \\
\bottomrule
\end{tabularx}
\caption{Joint set number \(J_n\).}
\label{tab:jn}
\end{table}

% ----------------------------------------------------------------------
\paragraph{Joint Roughness Number \(J_r\)}

\begin{table}[H]
\centering
\small
\begin{tabularx}{0.85\textwidth}{Y c}
\toprule
\textbf{Description} & \textbf{\(J_r\)} \\
\midrule
Discontinuous joints & 4 \\
Rough or irregular, undulating & 3 \\
Smooth, undulating & 2 \\
Slickensided, undulating & 1.5 \\
Rough or irregular, planar & 1.5 \\
Smooth, planar & 1 \\
Slickensided, planar & 0.5 \\
\bottomrule
\end{tabularx}
\caption{Joint roughness number \(J_r\).}
\label{tab:jr}
\end{table}

% ----------------------------------------------------------------------
\paragraph{Joint Alteration Number \(J_a\)}

\begin{table}[H]
\centering
\footnotesize
\begin{tabularx}{\textwidth}{Y c c}
\toprule
\textbf{Description} & \textbf{Residual friction angle} & \textbf{\(J_a\)} \\
\midrule
\multicolumn{3}{l}{\textbf{(a) Rock-wall contact, no mineral fillings, only coatings}} \\
Tightly healed, hard, non-softening, impermeable filling & -- & 0.75 \\
Unaltered joint walls, surface staining only & 25--35\(^{\circ}\) & 1 \\
Slightly altered, non-softening mineral coating, sandy particles & 25--30\(^{\circ}\) & 2 \\
Silty- or sandy-clay coatings, small clay fraction, non-softening & 20--25\(^{\circ}\) & 3 \\
Softening or low-friction mineral coatings & 8--16\(^{\circ}\) & 4 \\
\midrule
\multicolumn{3}{l}{\textbf{(b) Rock-wall contact before 10 cm shear, thin fillings}} \\
Sandy particles, clay-free disintegrated rock & 25--30\(^{\circ}\) & 4 \\
Strongly over-consolidated, non-softening clay fillings & 16--24\(^{\circ}\) & 6 \\
Medium or low over-consolidation, softening clay fillings & 12--16\(^{\circ}\) & 8 \\
Swelling clay fillings & 6--12\(^{\circ}\) & 8--12 \\
\midrule
\multicolumn{3}{l}{\textbf{(c) No rock-wall contact when sheared, thick fillings}} \\
Zones of disintegrated or crushed rock, strongly over-consolidated & 16--24\(^{\circ}\) & 6 \\
Zones of clay, medium or low over-consolidation & 12--16\(^{\circ}\) & 8 \\
Zones of clay, swelling clay & 6--12\(^{\circ}\) & 8--12 \\
Thick continuous zones of clay, strongly over-consolidated & 12--16\(^{\circ}\) & 10 \\
Thick continuous zones of clay, medium to low over-consolidation & 12--16\(^{\circ}\) & 13 \\
Thick continuous zones with swelling clay & 6--12\(^{\circ}\) & 13--20 \\
\bottomrule
\end{tabularx}
\caption{Joint alteration number \(J_a\).}
\label{tab:ja}
\end{table}

% ----------------------------------------------------------------------
\paragraph{Joint Water Reduction Factor \(J_w\)}

\begin{table}[H]
\centering
\small
\begin{tabularx}{\textwidth}{c Y c c}
\toprule
\textbf{Code} & \textbf{Conditions} & \textbf{Pressure (kg/cm\(^2\))} & \textbf{\(J_w\)} \\
\midrule
A & Dry excavations or minor inflow & \(<1\) & 1.0 \\
B & Medium inflow, occasional outwash of joint fillings & 1--2.5 & 0.66 \\
C & Jet inflow or high pressure in competent rock & 2.5--10 & 0.5 \\
D & Large inflow or high pressure, considerable outwash & 2.5--10 & 0.33 \\
E & Exceptionally high inflow, decaying with time & \(>10\) & 0.2--0.1 \\
F & Exceptionally high inflow, continuing without decay & \(>10\) & 0.1--0.05 \\
\bottomrule
\end{tabularx}
\caption{Joint water reduction factor \(J_w\).}
\label{tab:jw}
\end{table}

% ----------------------------------------------------------------------
\paragraph{Stress Reduction Factor (SRF)}
\label{sec:srf}

The SRF accounts for weakness zones, stress problems, squeezing, and swelling conditions.
It is one of the most judgement-sensitive parts of the Q-system.

\begin{table}[H]
\centering
\footnotesize
\begin{tabularx}{\textwidth}{Y c}
\toprule
\textbf{Description} & \textbf{SRF} \\
\midrule
\multicolumn{2}{l}{\textbf{(a) Weakness zones intersecting the excavation}} \\
Multiple weakness zones with clay, very loose surrounding rock & 10 \\
Multiple shear zones in competent clay-free rock & 7.5 \\
Single weak zone with or without clay, depth \(\leq 50\) m & 5.0 \\
Loose, open joints, heavily jointed or ``sugar cube'' rock mass & 5.0 \\
Single weak zone with or without clay, depth \(>50\) m & 2.5 \\
\midrule
\multicolumn{2}{l}{\textbf{(b) Competent rock, stress problems}} \\
Low stress, near surface, open joints, \(\sigma_{ci}/\sigma_1^{IS} > 200\) & 2.5 \\
Medium stress, favourable conditions, \(\sigma_{ci}/\sigma_1^{IS}=200\)--10 & 1.0 \\
High stress, very tight structure, \(\sigma_{ci}/\sigma_1^{IS}=10\)--5 & 0.5--2 \\
Moderate spalling/slabbing, \(\sigma_{ci}/\sigma_1^{IS}=5\)--3 & 5--50 \\
Spalling or rock burst, \(\sigma_{ci}/\sigma_1^{IS}=3\)--2 & 50--200 \\
Heavy rock burst, \(\sigma_{ci}/\sigma_1^{IS}<2\) & 200--400 \\
\midrule
\multicolumn{2}{l}{\textbf{(c) Squeezing rock}} \\
Mild squeezing, \(\sigma_{\theta}/\sigma_c = 1\)--5 & 5--10 \\
Heavy squeezing, \(\sigma_{\theta}/\sigma_c > 5\) & 10--20 \\
\midrule
\multicolumn{2}{l}{\textbf{(d) Swelling rock}} \\
Mild swelling & 5--10 \\
Heavy swelling & 10--15 \\
\bottomrule
\end{tabularx}
\caption{Stress Reduction Factor (SRF).}
\label{tab:srf}
\end{table}

\begin{warnbox}[Choosing the correct SRF category]
Category (b) applies when the rock mass is competent and the ratio
\(\sigma_{ci}/\sigma_1^{IS}\) can be meaningfully calculated.

Here:
\[
\sigma_1^{IS}=\text{maximum in-situ principal stress}
\]
not the Hoek--Brown failure stress \(\sigma_1^f\).

Category (a) applies when the excavation intersects weakness zones, heavily jointed zones,
or poorly interlocked rock, regardless of stress level.
\end{warnbox}

\begin{notebox}[Connection Chapter~\ref{ch:stress-intact-rock} notation]
In Chapter~\ref{ch:stress-intact-rock}, \(\sigma_1^{IS}\) was used for the maximum in-situ stress demand, while
\(\sigma_1^f\) was used for the major principal stress at failure from Hoek--Brown.
For SRF category (b), use \(\sigma_1^{IS}\), not \(\sigma_1^f\).
\end{notebox}

% ----------------------------------------------------------------------
\paragraph{Rock Mass Quality Based on Q}

\begin{table}[H]
\centering
\small
\begin{tabular}{c l}
\toprule
\textbf{Q-value} & \textbf{Rock mass quality} \\
\midrule
\(>400\) & Exceptionally good \\
100--400 & Extremely good \\
40--100 & Very good \\
10--40 & Good \\
4--10 & Fair \\
1--4 & Poor \\
0.1--1 & Very poor \\
0.01--0.1 & Extremely poor \\
\(<0.01\) & Exceptionally poor \\
\bottomrule
\end{tabular}
\caption{Rock mass quality based on Q-value.}
\label{tab:q_quality}
\end{table}

% ======================================================================
\section{Geological Strength Index}
\label{sec:geological-strength-index}

\subsection{Geological Strength Index (GSI)}
% ======================================================================

GSI describes rock mass structure and joint surface condition. It is used directly in
the generalised Hoek--Brown criterion Chapter~\ref{ch:rock-mass-parameters}.

\begin{keybox}[What GSI should represent]
GSI should reflect the \textbf{geological structure} and \textbf{joint surface condition} of
the rock mass. It should not directly include groundwater, excavation orientation, or
temporary construction effects.
\end{keybox}

% ----------------------------------------------------------------------
\paragraph{Structure Descriptions}

\begin{itemize}
 \item \textbf{Intact or massive}: intact rock specimens, few widely spaced discontinuities.
 \item \textbf{Blocky}: well interlocked, cubical blocks formed by approximately three joint sets.
 \item \textbf{Very blocky}: interlocked, partially disturbed, multifaceted angular blocks.
 \item \textbf{Blocky/disturbed/seamy}: folded or faulted rock with many intersecting discontinuities.
 \item \textbf{Disintegrated}: poorly interlocked, heavily broken mixture of rock fragments.
 \item \textbf{Laminated/sheared}: lack of blockiness due to weak foliation, schistosity, or shearing.
\end{itemize}

% ----------------------------------------------------------------------
\paragraph{Joint Surface Condition}

\begin{itemize}
 \item \textbf{Very good}: very rough, fresh, unweathered surfaces.
 \item \textbf{Good}: rough, slightly weathered, iron-stained surfaces.
 \item \textbf{Fair}: smooth, moderately weathered and altered surfaces.
 \item \textbf{Poor}: slickensided, highly weathered surfaces with compact coatings or fillings.
 \item \textbf{Very poor}: slickensided, highly weathered surfaces with soft clay coatings or fillings.
\end{itemize}

% ----------------------------------------------------------------------
\paragraph{Chart-Based GSI Estimation}

Chart-based GSI estimation uses two judgements:
\[
\text{Rock mass structure}
\quad + \quad
\text{Joint surface condition}
\quad \rightarrow \quad
\text{GSI range}
\]

\begin{notebox}[Recommended reporting style]
Do not report only one GSI number. State the descriptive basis, for example:

``The rock mass is interpreted as blocky to very blocky, with slightly rough but highly
weathered joint surfaces. Therefore, a GSI range of approximately XX--YY is considered
reasonable. A representative value of XX is adopted Chapter~\ref{ch:rock-mass-parameters} calculations.''
\end{notebox}

% ----------------------------------------------------------------------
\paragraph{Empirical GSI Correlations}

\subsubsection{GSI from RMR}

\[
\boxed{\text{GSI}=RMR_{\text{basic}}-5}
\]

where \(RMR_{\text{basic}}\) is the RMR calculated with:
\begin{itemize}
 \item groundwater rating set to 15, corresponding to completely dry conditions;
 \item orientation adjustment set to 0.
\end{itemize}

If the calculated RMR includes actual groundwater rating \(R_5\), then:
\[
RMR_{\text{basic}} = RMR - R_5 + 15
\]
and:
\[
\text{GSI} = RMR_{\text{basic}} - 5
\]

\begin{calcbox}[Example]
If:
\[
RMR = 60,\qquad R_5=7
\]
then:
\[
RMR_{\text{basic}} = 60 - 7 + 15 = 68
\]
\[
GSI = 68 - 5 = 63
\]
\end{calcbox}

\begin{warnbox}[Do not use raw RMR directly for GSI]
A common error is:
\[
GSI = RMR - 5
\]
without correcting the groundwater rating. This double-counts the groundwater penalty.
GSI should describe rock mass structure and joint surface condition, not groundwater.
\end{warnbox}

\subsubsection{Other Empirical Correlations}

Two additional empirical correlations are:
\[
GSI = 1.5J_{\text{cond89}} + \frac{RQD}{2}
\]
and:
\[
GSI =
\frac{52(J_r/J_a)}
{1+(J_r/J_a)}
+
\frac{RQD}{2}
\]

where \(J_{\text{cond89}}\) is the RMR joint condition rating \(R_4\), not the Q-system
joint roughness number \(J_r\).

\begin{warnbox}[Symbol caution]
\(J_{\text{cond89}}\) is the RMR joint condition rating. It is not the same as \(J_r\) in
the Q-system. Mixing these two symbols gives incorrect GSI values.
\end{warnbox}

% ======================================================================
\subsection{How to Select a Final GSI Value}
% ======================================================================

Engineers should compare:
\begin{itemize}
 \item chart-based GSI;
 \item GSI from \(RMR_{\text{basic}}-5\);
 \item GSI from \(1.5J_{\text{cond89}}+RQD/2\);
 \item GSI from the \(J_r/J_a\) correlation.
\end{itemize}

\begin{figure}[H]
\centering
\resizebox{0.95\textwidth}{!}{%
\begin{tikzpicture}[node distance=0.45cm, font=\small]

\node[workflowbox=orange] (chart)
{\textbf{Chart GSI}\\
\footnotesize structure + surface condition};

\node[workflowbox=blue, right=of chart] (corr)
{\textbf{Empirical GSI}\\
\footnotesize RMR, RQD, \(J_r/J_a\)};

\node[workflowbox=green, right=of corr] (compare)
{\textbf{Compare}\\
\footnotesize check consistency and outliers};

\node[workflowbox=purple, right=of compare] (select)
{\textbf{Adopt final GSI}\\
\footnotesize justify design value or range};

\draw[workflowarrow] (chart) -- (corr);
\draw[workflowarrow] (corr) -- (compare);
\draw[workflowarrow] (compare) -- (select);

\end{tikzpicture}
}
\caption{Workflow for selecting a defensible GSI value.}
\label{fig:gsi_selection}
\end{figure}

\begin{keybox}[Preferred approach]
Use the chart-based GSI as the primary geological judgement. Use empirical correlations
as checks. If the values differ significantly, explain why and adopt a conservative but
reasonable design value.
\end{keybox}

% ======================================================================
\section{Design data and calculation tables}
\label{sec:ch3-design-data}

\subsection{Project Input Data for design}
% ======================================================================

The following table summarises the key rock mass information for design.
Engineers must use the full project descriptions when assigning
ratings.

\begin{table}[H]
\centering
\scriptsize
\begin{tabularx}{\textwidth}{l Y Y Y Y}
\toprule
\textbf{Parameter} & \textbf{Granite} & \textbf{Sandstone} & \textbf{Mudstone} & \textbf{Limestone} \\
\midrule
\(\sigma_{ci}\) (MPa) & 265 & 110 & 27 & 85 \\
\(I_{s(50)}\) (MPa) & 12.2 & 5.3 & 2.3 & 5.0 \\
RQD (\%) & 70 & 68 & 38 & 52 \\
Joint condition & Tight, closed, very rough, undulating, unweathered & Slightly rough, highly weathered, not separated & Slickensided and planar, aperture 1--2 mm & Smooth, undulating, slightly weathered, clay-filled, aperture 1--2 mm \\
Joint sets / structure & Two joint sets from core observations & One sub-horizontal and one sub-vertical set & Poorly interlocked, heavily jointed, more than four sets & One joint set, spacing approximately 0.15 m \\
Groundwater & Dry & Medium inflow expected & Dripping at outcrop; medium inflow expected & Groundwater table 10 m below surface \\
Depth (m) & 780 & 450 & 100 & 80 \\
Stress ratio \(K\) & 3.2 & 1.8 & 1.0 & 2.0 \\
\bottomrule
\end{tabularx}
\caption{Summary of the design study rock mass information Chapter~\ref{ch:classification} classification.}
\label{tab:ch3-project-data}
\end{table}

\begin{notebox}[Do not treat this table as the full answer]
The table summarises key information only. Engineers still need to assign ratings,
show calculations, and justify decisions against the project description.
\end{notebox}

% ======================================================================
\subsection{Recommended Calculation Tables for Design Reports}
% ======================================================================

\paragraph{RMR Calculation Template}

\begin{table}[H]
\centering
\small
\begin{tabularx}{\textwidth}{l Y c Y}
\toprule
\textbf{RMR parameter} & \textbf{Description/value} & \textbf{Rating} & \textbf{Justification} \\
\midrule
\(R_1\): Intact strength & & & \\
\(R_2\): RQD & & & \\
\(R_3\): Joint spacing & & & \\
\(R_4\): Joint condition & & & \\
\(R_5\): Groundwater & & & \\
\midrule
\multicolumn{2}{r}{\textbf{Total RMR}} & & \\
\bottomrule
\end{tabularx}
\caption{Suggested RMR calculation table for each rock type.}
\label{tab:rmr_template}
\end{table}

\paragraph{Q-System Calculation Template}

\begin{table}[H]
\centering
\small
\begin{tabularx}{0.85\textwidth}{l Y c}
\toprule
\textbf{Parameter} & \textbf{Description/value} & \textbf{Adopted value} \\
\midrule
RQD & & \\
\(J_n\) & & \\
\(J_r\) & & \\
\(J_a\) & & \\
\(J_w\) & & \\
SRF & & \\
\midrule
\multicolumn{2}{r}{\textbf{Q-value}} & \\
\bottomrule
\end{tabularx}
\caption{Suggested Q-system calculation table for each rock type.}
\label{tab:q_template}
\end{table}

\paragraph{GSI Comparison Template}

\begin{table}[H]
\centering
\small
\begin{tabularx}{\textwidth}{l c Y}
\toprule
\textbf{Method} & \textbf{GSI value/range} & \textbf{Comment} \\
\midrule
Chart-based GSI & & Based on rock structure and surface condition \\
\(RMR_{\text{basic}}-5\) & & Groundwater corrected \\
\(1.5J_{\text{cond89}}+RQD/2\) & & Uses RMR joint condition rating \\
\(J_r/J_a\) correlation & & Uses Q-system roughness and alteration numbers \\
\midrule
\textbf{Adopted GSI} & & Justify final value or range \\
\bottomrule
\end{tabularx}
\caption{Suggested GSI comparison table for each rock type.}
\label{tab:gsi_template}
\end{table}

% ======================================================================
